\chapter{Metodologia e integração das componentes}
\label{chap:metodologia}

O protótipo Basketball AI Coach articula preparação de dados, aprendizagem
automática e planeamento. A descrição dos capítulos seguintes baseia-se nos
guias técnicos de \path{challenge-two/docs} e nos relatórios de experiências
conservados em \path{challenge-two/reports}. Os resultados apresentados
são os registados nesses artefactos; a elaboração deste relatório não
constitui uma nova execução das experiências.

\section{Organização do trabalho}

A divisão pelos seis elementos mantém três componentes em AAUTIA e três
em PLANTIA. O módulo de preparação de dados prepara os dados comuns, conforme o
Capítulo~\ref{chap:data}. O módulo de previsão de resultados prevê resultados de jogos
(Capítulo~\ref{chap:previsao-resultados}) e o módulo de avaliação de jogadores e alinhamentos avalia a produção
individual e as associações dos cinco iniciais com os resultados
(Capítulo~\ref{chap:desempenho}). Em PLANTIA, o módulo de planeamento de rotações seleciona o cinco
inicial e distribui minutos por programação por restrições
(Capítulo~\ref{chap:rotacoes}); o módulo de planeamento de substituições planeia substituições através
de um processo de decisão de Markov, comparando programação dinâmica e
algoritmo genético (Capítulo~\ref{chap:substituicoes}); e
o módulo de planeamento de vários jogos integra os módulos num horizonte de vários jogos
(Capítulo~\ref{chap:horizonte}). Cada componente abrange desenvolvimento,
testes, experiências e documentação.

\section{Questões de investigação e âmbito}

As experiências procuram responder às seguintes questões:
\begin{enumerate}
\item Qual dos modelos de classificação produz melhores probabilidades
 de vitória, face à frequência histórica de vitórias em casa e ao Elo?
\item Os modelos de regressão melhoram a previsão da produção individual
 e da margem final associada aos cinco iniciais, face a referências simples?
\item É possível transformar as estimativas disponíveis numa rotação
 completa que respeite posições, disponibilidade, minutos e descanso?
\item Como se compara uma política genética com a solução exata do MDP,
 e em que cenários estas políticas melhoram a utilidade simulada face
 à ordem fixa e a uma política gulosa?
\item Como compatibilizar planos de vários jogos com limites acumulados
 de minutos e quantificar a probabilidade de atingir um objetivo de vitórias?
\end{enumerate}

O âmbito implementado distingue três quantidades: probabilidades de
vitória, previsões estatísticas de produção e utilidades de planeamento.
Não existe uma relação causal estimada que converta uma alteração de
minutos ou uma substituição numa nova probabilidade de vitória. Assim,
a integração permite gerar planos e avaliar objetivos sob hipóteses
explícitas, mas não demonstra que esses planos aumentem as vitórias reais.

\section{Protocolo temporal comum}

A preparação documentada abrange 8\,283 jogos da fase regular da NBA, das
temporadas de 2019--20 a 2025--26. O treino utiliza 2019--20 a 2023--24,
a validação utiliza 2024--25 e o teste utiliza 2025--26. A cobertura é a
observada nas fontes, não uma certificação de completude do calendário.
Os subconjuntos efetivos diferem entre tarefas devido aos requisitos de
histórico individual e à qualidade dos registos de titulares.

As características históricas de cada dia são calculadas antes de
incorporar qualquer resultado desse dia. Imputação, normalização e modelos
são ajustados nos dados de treino; a seleção de configurações usa apenas
a validação. O modelo selecionado é reajustado com treino e validação,
permanecendo fixo durante o teste. Na avaliação diária, os históricos
podem incorporar jogos de dias de teste já concluídos.

O planeamento de vários jogos utiliza outro protocolo: congela as
observações antes de uma única data de corte para todo o horizonte.
Resultados intermédios ainda desconhecidos não são introduzidos no
histórico. Esta diferença é necessária para distinguir previsões diárias
de um plano elaborado antecipadamente para vários encontros.

\section{Interfaces, verificação e reprodutibilidade}

As tabelas preparadas ligam jogos, equipas e jogadores por identificadores
comuns. Os modelos guardam o pré-processamento e a informação necessária
à reprodução das previsões. Os planeadores verificam datas de treino,
compatibilidade das características e, nos modelos de desempenho,
as impressões digitais dos dados de entrada.

A sequência de execução é preparação, treino dos módulos de previsão de resultados e de avaliação de jogadores e alinhamentos e
planeamento. O módulo de planeamento de rotações consome os modelos do módulo de avaliação de jogadores e alinhamentos; o módulo de planeamento de substituições
consome a rotação do módulo de planeamento de rotações. O módulo de planeamento de vários jogos reutiliza diretamente as
funções de previsão, otimização e MDP, sem exigir relatórios autónomos dos
módulos de planeamento de rotações e de substituições e sem voltar a treinar os modelos durante o planeamento.

Os guias documentam testes de exclusão de informação futura, validação
das soluções e consistência da integração. As configurações, métricas,
identificadores e hashes permitem relacionar cada resultado com as suas
entradas. O marcador \texttt{\_SUCCESS} identifica execuções concluídas;
um diretório com ficheiros parciais não constitui uma experiência válida.
